% !TEX root = ../main.tex
\PlanPage{对照与消融：识别每个结构化模块的贡献}{吕洪宇}{31}
\noindent\textbf{完整工作流的价值应通过同一任务、同一材料和同一复核标准来检验，而不是由系统自身的叙述定义。}因此，项目将人工原始材料检索、通用摘要/问答、文档检索、结构化完整流程及两个关键模块移除设置纳入预注册对照。比较的重点是路径是否更准确、证据是否更可定位、主体是否更一致、完成研究任务所需的时间与返工是否减少；每个指标均对应清晰的判定对象和原始记录。

\begin{center}\small 表31-1\quad 对照与消融设置的研究问题设计\par\vspace{1mm}
\begin{tabularx}{\linewidth}{p{0.20\linewidth}p{0.33\linewidth}X}
\toprule
设置 & 保留或移除的研究能力 & 主要回答的问题\\
\midrule
人工原始材料检索 & 研究者自行阅读公开材料、形成路径和公司清单 & 在无系统帮助时，完成同一任务需要何种时间和人工整理成本？\\
通用摘要/问答 & 输入相同问题，但不要求显式路径、证券主体和证据ID & 缺少结构化约束时，事实、主体和引用是否容易发生漂移？\\
文档检索流程 & 返回相关片段和链接，不执行图谱查询与评分编排 & 仅提高材料可得性，能否形成可审阅的产业传导和公司解释？\\
完整工作流 & 事件、路径、公司、评分和证据经总控组合为报告 & 跨模块结构化编排是否提升路径、引用和队列的一致性？\\
去图谱路径 & 保留证据和评分，但隐藏路径与规则约束 & 产业路径是否提升研究者解释传导机制和复核公司关联的能力？\\
去证据校验 & 保留路径与排序，但取消\texttt{evidence\_ids}强制绑定 & 证据约束是否提高结论可定位率与事实支持率？此设置仅用于离线实验。\\
\bottomrule
\end{tabularx}\end{center}

所有设置使用同一事件题、相同截止日期的材料、相同公司范围和相同答题模板。至少两名复核者在不知晓方法来源的条件下，对事件事实、路径边、公司主体与证据定位分别判定；项目记录路径边精确率与召回率、证据可定位率、事实支持率、主体匹配率、评分复算一致率及完成时间中位数。对于路径解释类任务，除数值指标外，复核表还记录错误发生在事件识别、关系方向、公司映射、财务解释还是证据引用层。

这样设计消融实验，能够让不同模块承担明确的证据角色：图谱模块检验传导解释，证据模块检验引用完整性，评分模块检验计算可复现性，总控模块检验跨模块一致性。正式实验将在数据主链接入后按预注册协议运行；在此之前，项目以已实现的契约测试、接口验证和运行记录作为工程层证据，为后续研究评测建立稳定基线。
